\begin{papercard}{bfarm2023}
\cf{Stages} \stagetag{S} (what counts as a DiGA and so what the sample frame is), \stagetag{A} (eligibility and sampling frame via the DiGA directory), \stagetag{C} (legal expectations on consent, tracking and advertising that a policy can be coded against), \stagetag{R} (public information channels for patients). BfArM, ``The Fast-Track Process for Digital Health Applications (DiGA) according to Section 139e SGB V: A Guide for Manufacturers, Service Providers and Users'' (English version; 128 PDF pages, the document's own numbering runs to 126) \pc{bfarm2023}{1}{The Fast-Track Process for Digital Health Applications (DiGA) according to Section 139e SGB V}. It is an official guide to a German regulatory procedure, not research. The footnote says the English text is a service and that ``the corresponding standards in the German version are legally binding'' \pc{bfarm2023}{12}{The corresponding standards in the German version are legally binding}. The file states no version date; the wording is that of 2020 (deadlines of 2021 are future tense, the 16 July 2020 judgment on Privacy Shield is the latest event cited), although the corpus label says 2023, so check the current edition before citing details. I read chapters 1--3 and the parts of 4 and 5 on evidence and life cycle; the evidence-study design sections (4.3--4.6), fees and the glossary were not read in detail.

\cf{The fast track} The Digital Healthcare Act came into force on 19 December 2019 and introduced the ``app on prescription'': statutory-insured people (about 73 million) may receive prescribed DiGA, reimbursed by health insurers \pc{bfarm2023}{8}{came into effect on December 19th 2019 introducing the app on prescription as part of healthcare provided to patients}. The BfArM has three months from a complete application to assess the manufacturer's statements ``from data protection to user friendliness'' and the evidence of a positive healthcare effect \pc{bfarm2023}{8}{The essence of this assessment is the examination of the manufacturer’s statements about the product qualities – from data protection to user friendliness - and the examination of the evidence of the positive healthcare effect}. Without sufficient evidence a \emph{provisional} listing with a one-year (exceptionally two-year) trial is possible \pc{bfarm2023}{9}{If the manufacturer cannot provide sufficient evidence for a positive healthcare effect but all other requirements are fulfilled, it will be possible to apply for a provisional listing}. The idea: link ``privacy and information security on the one hand and user friendliness and high performance on the other'' \pc{bfarm2023}{10}{It aims for a successful link between privacy and information security on the one hand and user friendliness and high performance on the other hand}, with ``maximum transparency'' and a public directory \pc{bfarm2023}{10}{Maximum transparency is a key focus}.

\cf{What is a DiGA} A medical device of risk class I or IIa whose main function is digital, supporting recognition, monitoring, treatment or alleviation of disease, not primary prevention, used by the patient or by patient and provider \pc{bfarm2023}{13}{A DiGA is a medical device that has the following properties}; ``digital assistants in the hands of patients'' \pc{bfarm2023}{13}{DiGA are therefore digital assistants in the hands of patients}. Apps that only collect sensor data without patient interaction are not DiGA. An app earlier offered in a store with an advertising-based model can be listed; the BfArM examines only the version or variant applied for and does not look at parallel versions \pc{bfarm2023}{14}{That is permissible. The BfArM only examines the version or variant of a DiGA for which an application for inclusion in the DiGA directory is made}. The 2.1 sub-sections give examples of what is and is not a DiGA (hardware and services combinations, scope, prevention).

\cf{The DiGA directory (public channel)} Its main goal is ``to enable and strengthen a trusting use of DiGA'' \pc{bfarm2023}{24}{The most important goal of the DiGA directory is to enable and strengthen a trusting use of DiGA in healthcare by patients and healthcare providers}. Information shown to insured people includes a plain description of goal, principle and use, functions, the \emph{confirmed Annex 1 checklist on data protection and information security} and Annex 2 quality checklist, extra costs such as in-app purchases, and \emph{locations where the data is processed} \pc{bfarm2023}{25}{Confirmed checklist regarding data protection and information security (Annex 1 DiGAV)}; \pc{bfarm2023}{25}{Locations where the data is processed}. Contact and privacy-question addresses are not listed in the directory because duties to inform sit in other law such as Article 13 GDPR \pc{bfarm2023}{26}{these obligations to provide information are regulated in other laws and regulations which are applicable}. Providers see evidence status, price and intended user roles; professionals see studies and study reports. Technical information includes the manufacturer's confirmation of compatibility \pc{bfarm2023}{27}{Confirmation of compatibility by the manufacturer regarding supported platforms, devices and potential additional products}. Directory data are to be available to third parties by an API from 2021 \pc{bfarm2023}{28}{will be made available to other interested public and non-profit institutions via an Application Programming Interface from 2021}, and manufacturers may mark information as confidential because of trade secrets, which they must nevertheless enter in the application \pc{bfarm2023}{28}{The manufacturer can mark the information in the application procedure which may not be made public due to legal requirements}.

\cf{Checklists and what the BfArM verifies} Manufacturers prove requirements mainly through the completed yes/no checklists of Annexes 1 (data protection and information security) and 2 (interoperability, robustness, consumer protection, usability, provider support, medical quality, patient safety) \pc{bfarm2023}{36}{Manufacturers have to demonstrate this to the BfArM with emphasis on the completed checklists of the appendices 1 and 2 of the DiGAV}; the BfArM may request further evidence and must receive free login access to the DiGA \pc{bfarm2023}{36}{In any case, free access (login data) to the DiGA must be provided to the BfArM}. Compliance is given when all statements are ticked ``applicable'' \pc{bfarm2023}{36}{All statements are formulated in a way that the compliance with the required criteria is given when all associated statements have been ticked as applicable}; a ``not applicable'' not pre-specified needs a written justification, otherwise the application is refused \pc{bfarm2023}{37}{A not applicable answer not specified for selection in Annex 1 or 2 requires a written justification why the overall criteria of the statement are nevertheless fulfilled}. For safety the BfArM ``only carries out checks on the formal legality of the CE marking'' \pc{bfarm2023}{37}{As a rule, the BfArM only carries out checks on the formal legality of the CE marking for this requirement}. Annex 1 has \textbf{40 statements} covering technical and organisational measures and processes \pc{bfarm2023}{39}{The checklist to be completed by the manufacturer of the DiGA for the application contains 40 statements}. In the chapters I read, verification is documentary (checklists, CE marking, login access); I found no description of testing the statements against the app's network behaviour.

\cf{Data protection rules specific to DiGA} Users must be able to rely on lawful, careful handling of their data \pc{bfarm2023}{38}{Insured persons as users of a DiGA must be able to rely on the fact that the manufacturer complies with legal requirements for data protection}; health data are mainly regulated by Section 22 BDSG with Article 9 GDPR \pc{bfarm2023}{38}{The main regulation for the processing of health data is Section 22 BDSG (in conjunction with Article 9 GDPR, if applicable)}. Section 4(2) DiGAV \emph{restricts consent} to four purposes: (1) intended use, (2) proof of positive healthcare effect in a trial, (3) verification of performance-related price agreements, (4) maintenance, usability and further development. Consent for other purposes is not allowed \pc{bfarm2023}{39}{It is not permitted to obtain consent of the person concerned to legitimise the processing of health data for purposes other than those specified in Section 4 paragraph 2 clause 1 DiGAV}. Consent must be explicit, informed and voluntary at the start \pc{bfarm2023}{40}{At the beginning of the use of DiGA and prior to the collection and further processing of personal data by the DiGA, a voluntary and informed, explicit consent of the person concerned must be obtained}; GDPR data minimisation and privacy by design and default ``continue to apply without restriction'' \pc{bfarm2023}{40}{The standards of the GDPR, in particular data minimisation and data protection-friendly technology design through Privacy by Design and Privacy by Default, continue to apply without restriction}. A ``record everything and see later'' approach is ``not according to the guidelines'' \pc{bfarm2023}{41}{An approach following the principle we record everything we can get and then we see what we really need is not according to the guidelines}. Purposes 1--3 may be combined in one consent, but optional parts must be recognisable and refusable and the coupling prohibition (Article 7(4) GDPR) applies \pc{bfarm2023}{41}{may not be made dependent on consent to the processing of personal data which is not necessary for the fulfilment of the contract}. Purpose 4 consent is optional: refusal must not affect use \pc{bfarm2023}{42}{The insured person must be able to continue using the DiGA without restrictions}, and ``comprehensive tracking of user activities is not permitted'' \pc{bfarm2023}{42}{However, comprehensive tracking of user activities is not permitted}. ``A general coupling prohibition applies to the DiGA'': paying with data for features is not permitted \pc{bfarm2023}{43}{A general coupling prohibition applies to the DiGA, i.e. a payment with data for features within a DiGA is not permitted}. Processing outside the EU/EEA/Switzerland is limited to countries with an Article 45 adequacy decision; standard contractual clauses and binding corporate rules are not allowed for DiGA \pc{bfarm2023}{44}{Processing of personal data outside the EU on the basis of Article 46 GDPR (standard contractual clauses) or Article 47 (Corporate Binding Rules) is not permitted for DiGA}, and the guide concludes that processing of health data in the USA ``is therefore not permissible'' \pc{bfarm2023}{44}{Processing of health data in the USA is therefore not permissible for a DiGA} (a conclusion tied to the legal situation of 2020; check later developments before relying on it).

\cf{Information security} Requirements cover confidentiality, integrity and availability \pc{bfarm2023}{45}{The requirements concerning data protection refer to the protection of confidentiality, integrity and availability of all data processed via a DiGA} and follow BSI standards and IT-Grundschutz \pc{bfarm2023}{45}{The specifications in the DiGAV for information security are based on the relevant publications and recommendations of the Federal Office for Information Security (BSI)}; additional requirements apply for very high protection needs; interoperability (export in standard and human-readable formats, device interfaces) and robustness follow in 3.5--3.6.

\cf{Consumer protection} The starting point is fairness: users are in ``a special life and/or illness situation'' that ``must not be exploited by the manufacturer'' \pc{bfarm2023}{67}{users of DiGA find themselves in a special life and / or illness situation simply because of their motivation to use a particular DiGA, which must not be exploited by the manufacturer}. Manufacturers must be transparent about purpose and functionality and show on the sales platform or website which features cost what; the guide admits store operators severely restrict what can be displayed \pc{bfarm2023}{67}{in the case of apps for iOS or Android, the scope of what must or may be displayed in the App Store and to what extent is usually severely restricted by the App Store operator}. Users must be able to judge compatibility before installing \pc{bfarm2023}{68}{Before installing a DiGA or using it for the first time, insured persons must be able to determine to what extent the application fits their own requirements}. In-app purchases: not advertised in the DiGA, no auto-renewing subscriptions or time-limited offers, and accidental purchase must be excluded \pc{bfarm2023}{69}{In-app purchases shall not be automatically renewable subscriptions or limited-time special offers}. \textbf{Advertising is prohibited}: ``digital health applications must be free of advertising'', own or third-party \pc{bfarm2023}{69}{does not differentiate whether it is the manufacturer's own advertising for his own products or third-party advertising for third-party offers}, including addressing users ``on an emotional level with the aim of bringing about an irrational decision'' \pc{bfarm2023}{69}{This also includes practices of courting or addressing on an emotional level with the aim of bringing about an irrational decision}. Direct marketing by other channels would need consent, which is not allowed for such purposes \pc{bfarm2023}{70}{No. This would require user consent. However, consent may only be obtained for the purposes specified}.

\cf{Evidence of benefit (chapter 4, overview)} A positive healthcare effect is a medical benefit or a patient-relevant improvement of structure and processes \pc{bfarm2023}{77}{positive healthcare effects (positive Versorgungseffekte, pVE) are either a medical benefit}; effects must be patient-centred, and staff workload or economic indicators do not count \pc{bfarm2023}{77}{The workload of medical staff or economic indicators of healthcare are not patient-relevant endpoints}; medical benefit means patient-relevant endpoints such as morbidity, mortality or quality of life \pc{bfarm2023}{78}{Those who claim a medical benefit for a DiGA must show that patient-relevant endpoints, in particular morbidity, mortality or quality of life, are positively influenced}. Harms and data risks are not part of the evidence definition; privacy sits in Annex 1.

\cf{Life cycle and changes} Manufacturers must check each change against ``significant change'' (Section 18 DiGAV) and report it \pc{bfarm2023}{109}{Every time a manufacturer makes changes to the DiGA he has to check whether these are significant changes according to Section 18 DiGAV}; a change in the location of data storage or in prices counts as significant, new hardware or OS support does not \pc{bfarm2023}{112}{A significant change can be, for example, a change in the location of data storage, a change in prices, etc}. Unnotified changes (due three months before) can lead to deletion \pc{bfarm2023}{112}{If the BfArM becomes aware of changes made which have not been notified to the BfArM or have not been notified in time (three months before the change is made), the BfArM may delete the affected DiGA}. Information linked from the directory (store page, website) must be kept current by the manufacturer ``in an active and independent manner'' \pc{bfarm2023}{110}{All information linked in the DiGA directory in the distribution platform or on the application website must be kept up to date by the manufacturer in an active and independent manner}; Annex 1 also implies deletion of data no longer needed \pc{bfarm2023}{110}{deleting or blocking of data that no longer required}.

\cf{Limits (card)} A regulator's guide for manufacturers: it states rules and interpretations, not evidence on whether DiGA comply or whether users understand the information; verification is largely documentary (checklists, CE marking, login access); the statements on third countries and the USA reflect 2020 law; version date unstated in the file; only low-risk devices (class I/IIa) with a medical purpose are in scope, so DiGA are a small, regulated subset of the health and mental-health apps in the stores; the English text is not the legally binding one.

\cf{Use for this thesis} (1) \emph{Sampling frame and definition}: DiGA are a clearly bounded, publicly listed population with a stated definition of scope; if the thesis includes DiGA it can cite this guide for the eligibility definition and use the directory as a source for the manufacturer, listing status and data-processing locations; DiGA versus general store apps is a natural comparison for stage A. (2) \emph{Expectations to code against}: for DiGA the guide implies that a policy should not seek consent for marketing or general analytics, should not tie features to data, should not claim processing in non-adequate third countries, and the app should contain no advertising; deviations are findings about communication or practice, to be phrased as inconsistency with the guide's stated rules, not as legal violations. (3) \emph{Communication channels}: the directory publishes the confirmed Annex 1 checklist and processing locations for patients; the thesis can ask whether the policy text and store label agree with those directory entries and whether the checklists are accessible to a layperson (the guide explicitly says contact and privacy-question details are not in the directory). (4) \emph{Harm framing}: the guide's fairness and ``exploiting a special illness situation'' language is a regulator-side harm statement that supports a harm-tree node on exploitation of vulnerability and on manipulative advertising. (5) \emph{Change obligations}: a policy change on data location is a ``significant change'' that must be notified, which supports capturing policy versions and dates (stage D). (6) Use with the WP29 transparency guidelines (what the notice must contain) and with the GDPR; cite the German original for legal claims.
\end{papercard}
